% ============================================================================
\section{Synthesis and discussion}
% ============================================================================

\begin{frame}[t]{Summary: Transformers 2017 vs.\ 2026}
  \footnotesize
  \vspace{-0.2em}
  \begin{center}
    \renewcommand{\arraystretch}{1.25}
    \begin{tabular}{@{}>{\raggedright\arraybackslash}p{0.20\linewidth}>{\raggedright\arraybackslash}p{0.37\linewidth}>{\raggedright\arraybackslash}p{0.37\linewidth}@{}}
      \textbf{Component} & \textbf{2017 Transformer} & \textbf{2026 consensus} \\
      \hline
      Attention mechanism & Multi-Head Attention (MHA) & Grouped-Query Attention (GQA) / Multi-Head Latent Attention (MLA) \\
      \hline
      Position encoding & Absolute sinusoidal (or learned absolute) & Rotary Positional Embeddings (RoPE) \\
      \hline
      Normalization & Post-Layer Normalization (post-LN) & Pre-Layer Normalization with RMSNorm \\
      \hline
      Activation & ReLU in the FFN & SwiGLU / gated linear units \\
      \hline
      Bias terms & Biases in linear layers common & Bias-free or minimal-bias architectures \\
      \hline
      Parameter density & Dense (all parameters active per forward) & Sparse activation (e.g., Mixture-of-Experts) \\
      \hline
    \end{tabular}
  \end{center}
\end{frame}

\begin{frame}{Summary}
  \begin{itemize}
    \item The slowest part is often getting data in and out of memory, not just the math.
    \item New models (2024–2026) are designed together, optimizing attention, normalization, position encoding, sparsity (MoE), and serving (memory, paging, kernels).
    \item The main goals: cut down on moving data, adjust compute to the input’s length and type, and add more parameters without needing a lot more computation per token.
  \end{itemize}
\end{frame}
